\documentclass[11pt]{article}

\usepackage[margin=1in]{geometry}
\usepackage{amsmath,amssymb}
\usepackage{booktabs}
\usepackage[american]{circuitikz}
\usepackage{xcolor}
\usepackage{enumitem}

\newenvironment{solution}
  {\par\medskip\noindent\textbf{Solution.}\ }
  {\par\medskip}

\begin{document}

\section*{Problem 4 \hfill (13 points total)}

A black box filled with linear circuitry is connected to a variable resistor.
An ideal ammeter (with zero resistance) and an ideal voltmeter (with infinite
resistance) are used to measure current and voltage as shown below. The results
of several sample readings are listed in the accompanying table.

\begin{center}
\begin{minipage}[c]{0.55\textwidth}
\centering
\begin{circuitikz}[scale=1.0, transform shape]
  % Black box
  \draw[fill=blue!8] (0,0) rectangle (2.2,3.6);
  \node[align=center] at (1.1,2.75) {Black box\\(linear network)};

  % Output terminals and load loop
  \draw (2.2,2.5) -- (3.2,2.5)
        to[ammeter, l=$I$] (4.4,2.5)
        -- (5.25,2.5)
        to[vR, l=$R$] (5.25,0.6)
        -- (2.2,0.6);

  % Voltmeter across output
  \draw (3.2,2.5) to[voltmeter, l_=$V$] (3.2,0.6);
\end{circuitikz}
\end{minipage}%
\begin{minipage}[c]{0.38\textwidth}
\centering
\begin{tabular}{ccc}
\toprule
$R\;(\Omega)$ & $V\;(\mathrm{V})$ & $I\;(\mathrm{A})$ \\
\midrule
4 & 6.0 & 1.5 \\
6 & 8.0 & 1.333 \\
8 & 9.6 & 1.2 \\
\bottomrule
\end{tabular}
\end{minipage}
\end{center}

\begin{enumerate}[label=(\alph*), leftmargin=2em]
  \item Find $I$ and $V$ when $R=17\,\Omega$. \hfill (7 points)

  \begin{solution}
  Replace the black box by its Th\'evenin equivalent, consisting of a source
  $V_{\mathrm{th}}$ in series with a resistance $R_{\mathrm{th}}$. For a load
  resistance $R$,
  \[
    \frac{V}{V_{\mathrm{th}}}=\frac{R}{R+R_{\mathrm{th}}}.
  \]
  Using the measurements for $R=4\,\Omega$ and $R=6\,\Omega$ gives
  \[
    \frac{6}{V_{\mathrm{th}}}=\frac{4}{4+R_{\mathrm{th}}},
    \qquad
    \frac{8}{V_{\mathrm{th}}}=\frac{6}{6+R_{\mathrm{th}}}.
  \]
  Eliminating $V_{\mathrm{th}}$,
  \begin{align*}
    \frac{6}{4}(4+R_{\mathrm{th}})
      &= \frac{8}{6}(6+R_{\mathrm{th}}),\\
    6+\frac{3}{2}R_{\mathrm{th}}
      &= 8+\frac{4}{3}R_{\mathrm{th}},\\
    R_{\mathrm{th}} &= 12\,\Omega.
  \end{align*}
  It follows that
  \[
    V_{\mathrm{th}}
      =\frac{6}{4}(4+12)
      =24\,\mathrm{V}.
  \]
  Therefore, with $R=17\,\Omega$,
  \[
    V=V_{\mathrm{th}}\frac{R}{R+R_{\mathrm{th}}}
      =24\frac{17}{17+12}
      =\frac{408}{29}\,\mathrm{V}
      \approx 14.1\,\mathrm{V},
  \]
  and
  \[
    I=\frac{V}{R}
      =\frac{24}{29}\,\mathrm{A}
      \approx 0.83\,\mathrm{A}.
  \]
  \end{solution}

  \item Determine the value of $R$ that achieves maximum power transfer from
  the black box to the variable resistor. \hfill (2 points)

  \begin{solution}
  Maximum power transfer occurs when the load resistance equals the Th\'evenin
  resistance:
  \[
    \boxed{R=R_{\mathrm{th}}=12\,\Omega.}
  \]
  \end{solution}

  \item What is the maximum power transferred to the variable resistor?
  \hfill (4 points)

  \begin{solution}
  With $R=R_{\mathrm{th}}$, the maximum load power is
  \[
    P_{\max}=\frac{V_{\mathrm{th}}^2}{4R_{\mathrm{th}}}
      =\frac{(24\,\mathrm{V})^2}{4(12\,\Omega)}
      =\boxed{12\,\mathrm{W}}.
  \]
  \end{solution}
\end{enumerate}

\end{document}
